\section{Methodology}\label{s:method}

We study a two-layer attention-only transformer trained on the ordered bigram copying task introduced by \citet{farrell2025order} (Table \ref{tab:model_config}). The model encodes bigram $(d_1,d_2)$ into a single separator token, $s$, and reconstructs it as $(o_1,o_2)=(d_1,d_2)$. The 5-token input is represented as $[d_1, d_2, s, o_1, o_2]$, where $d_1$ and $d_2$ are uniformly sampled from $[0,99]$ and $s, o_1, o_2$ are special tokens. This gives 10,000 possible ordered bigrams.
\citet{farrell2025order} find that the residual at $s$ after the first layer, $\boldsymbol r_{s}^{L_1}$, is:
\begin{equation*}
    \boldsymbol r^{L_1}_s = \alpha_{s\to d_1} (\boldsymbol E_{d_1} + \boldsymbol P_{d_1}) + \alpha_{s \to d_2}(\boldsymbol E_{d_2} + \boldsymbol P_{d_2}) + \boldsymbol E_s + \boldsymbol P_s
\end{equation*}
where $\alpha_{s\to d_i}$ is the attention weight in the first layer for query $s$ (separator token) and key $d_i$ (element of input bigram), and $\boldsymbol E_{x}$ and $\boldsymbol P_{x}$ are the token and positional embeddings for the token at $x$ respectively.
Thus, the model represents bigram orderings at $s$ using a linear combination of basis vectors, $(\boldsymbol E_{b} + \boldsymbol P_{d_2})$ and $(\boldsymbol E_{b} + \boldsymbol P_{d_2})$, with coefficients defined by the attention probabilities.
% the relative magnitude of scalar attention weights, rather than representing each possible ordering as a different direction in activation space (which other composition methods do \citep{wattenberg2024relational}).


Models can represent more than $n$ features in $n$-dimensional activations through superposition, which complicates interpretability \citep{elhage2022toy}. SAEs decompose dense, superposed model activations into sparse feature representations \cite{bricken2023monosemanticity}.
We train BatchTopK \cite{bussmann2024batchtopk}, JumpReLU \cite{rajamanoharan2024jumping} and Matryoshka \cite{bussmann2025learning} SAEs on $\boldsymbol r_{s}^{L_1}$ using implementations from \citet{marks2024dictionary_learning}. 
SAEs are trained on residual activations on the full finite input space of the toy task (10,000 inputs). 
Because our goal is not out-of-distribution generalisation, we do not need a train/test split. Instead, we evaluate SAEs on sparsity ($L_0$) and loss recovered \cite{karvonen2025saebench,ferrando2024primer} over the complete activation distribution.

\begin{table}[h]
\caption{Selected SAE configuration. The optimiser, betas, schedule and auxiliary loss are fixed in the implementation \cite{marks2024dictionary_learning} and excluded here for brevity.}
\label{tab:sae_config}
\centering
\setlength{\tabcolsep}{2pt}
\begin{tabular}{ll}
\toprule
\textbf{Parameter} & \textbf{Value} \\
\midrule
SAE type & BatchTopK \\
% Activation source & SEP token, residual stream after layer 0 \\
Input dimension ($d_{\text{model}}$) & 64 \\
Dictionary size ($d_{\text{SAE}}$) & 128 \\
$k$ & 3 \\
Actual sparsity ($L_0$) & 3.00 \\
Activation centring & Mean-centred (pre-training) \\
Encoder init & $W_{\text{enc}} = W_{\text{dec}}^\top$ \\
Decoder normalisation & Unit-norm columns \\
\midrule
Learning rate & $3 \times 10^{-4}$ \\
% Optimiser & \texttt{Adam} \\
% Betas & $(0.9,\ 0.999)$ \\
Batch size & 4096 \\
Training steps & 50{,}000 \\
% LR schedule & Linear warmup (1{,}000 steps), then constant \\
% Auxiliary loss ($\alpha$) & $1/32$ \\
Seed & 1 \\
\midrule
Hardware & NVIDIA GeForce RTX 2080 Ti \\
\bottomrule
\end{tabular}
\end{table}

After a grid search (Table \ref{tab:sae_sweep}) we find that BatchTopK SAEs with $k \in [3,5]$ all achieve similar top performance, so we select $k=3$. We select $d_\text{SAE} = 128$ which achieves similar performance to larger $d_\text{SAE}$. We filter for loss recovered over 0.999 and use the checkpoint with the fewest dead latents ($3.1\%$). Table \ref{tab:sae_config} shows the final configuration.
% This checkpoint selection introduces selection bias toward high-performing models. We partially address this by checking whether special latents (identified by $\Delta\alpha$ correlation) appear across the broader SAE sweep, while reserving steering for selected high-performing checkpoints.

We inspect the activating examples for each latent, as well as their firing rates and correlation with the attention weights $\alpha$ from the first layer.
Moreover, \citet{brien2025steering} steer model activations at inference time by amplifying SAE latent activations mediating refusal. 
We apply this principle to test if a latent has graded activations, expecting that the model output will change at different activation graduations. We linearly scale the SAE latent activation for a given input. We then patch the SAE reconstruction back into the model downstream and compare the new output to the original output. 
We use the checkpoints in Table \ref{tab:sae_sweep_configs} to test whether special latents appear across architectures and hyperparameters. The correlation analysis in Figure \ref{fig:special_latents_vs_performance} covers all of our trained SAEs. Because steering is much more expensive, we run it on six selected high-performing checkpoints (Table \ref{tab:selected-saes}) chosen to span architecture, dictionary size and sparsity settings. Though we note this selected set is insufficient for statistical significance.
